\documentclass[12pt]{article}
\usepackage[utf8]{inputenc}
\usepackage[spanish]{babel}
\usepackage{geometry}
\usepackage{hyperref}
\usepackage{booktabs}
\usepackage{longtable}

\geometry{a4paper, margin=2.5cm}

\title{Memoria Técnica y Estratégica: Análisis Operativo, Arquitectura y Potencial de Implementación de la Red Algorand (ALGO)}
\author{Imad Habib Jali (901421) \and Jorge Bellido Lobera (903080)}
\date{}

\begin{document}

\maketitle

\section*{Resumen Ejecutivo}
\textbf{A:} Junta Directiva / Comité de Inversiones Estratégicas \\
\textbf{De:} Analista Experto en Tecnología Blockchain \\
\textbf{Asunto:} Análisis Operativo, Arquitectura y Potencial de Implementación de la Red Algorand (ALGO)

\noindent\rule{\textwidth}{0.4pt}

\section{¿Qué es Algorand?}

Algorand es una infraestructura blockchain de tercera generación, pública y descentralizada, concebida para resolver el histórico "Trilema Blockchain": el desafío de lograr escalabilidad, seguridad y descentralización de forma simultánea y sin concesiones. 

Fundada en 2017 por Silvio Micali ---criptógrafo, profesor del MIT y galardonado con el Premio Turing---, la red principal (Mainnet) se lanzó en 2019. A diferencia de las redes de primera generación concebidas como simples sistemas de pago, Algorand actúa como un ecosistema integral para la economía digital moderna. Permite la ejecución de contratos inteligentes (Smart Contracts), el despliegue de aplicaciones descentralizadas (dApps) y la tokenización de activos de grado institucional. 

Destaca en el panorama tecnológico por su enfoque en la adopción empresarial y gubernamental, siendo una de las plataformas preferidas para el desarrollo de Monedas Digitales de Bancos Centrales (CBDC), como la de las Islas Marshall, y por mantener un compromiso de huella de carbono negativa.

\section{¿Cómo funciona esta red blockchain?}

La superioridad técnica de Algorand se fundamenta en dos innovaciones arquitectónicas principales: su mecanismo de consenso y su estructura funcional separada en capas.

\subsection{Mecanismo de Consenso: Pure Proof of Stake (PPoS)}
A diferencia de los sistemas tradicionales, el protocolo PPoS no requiere minería masiva ni el bloqueo forzoso de grandes capitales. Su funcionamiento se basa en la selección criptográfica y aleatoria:
\begin{itemize}
    \item \textbf{Selección Secreta y Aleatoria:} Mediante una Función Aleatoria Verificable (VRF), el sistema selecciona en microsegundos a un líder para proponer un bloque y a un comité para validarlo.
    \item \textbf{Seguridad por Opacidad:} Nadie sabe quién es el validador hasta que el bloque ya ha sido propuesto. Esto hace que la red sea matemáticamente invulnerable a ataques de denegación de servicio (DDoS) o sobornos a los nodos, ya que los atacantes no saben a quién atacar.
    \item \textbf{Recompensas Inclusivas:} Cualquier usuario que posea al menos 1 ALGO en su billetera es elegible para participar en el consenso y la gobernanza, democratizando la seguridad de la red.
\end{itemize}

\subsection{Arquitectura de Doble Capa}
Para evitar la congestión que paraliza a otras redes, Algorand divide la carga computacional:
\begin{itemize}
    \item \textbf{Capa 1 (Layer 1):} Procesa transacciones cotidianas, transferencias de tokens nativos (mediante el protocolo ASA) y contratos inteligentes simples a velocidades ultra rápidas (miles de transacciones por segundo con confirmación final en menos de 5 segundos).
    \item \textbf{Capa 2 (Layer 2):} Se reserva para contratos inteligentes de alta complejidad computacional. Estas operaciones se ejecutan fuera de la cadena principal y, periódicamente, envían certificados de validez a la Capa 1, manteniendo la red base ligera, rápida y económica.
\end{itemize}

\section{Diferencias estructurales con Bitcoin y Ethereum}

El posicionamiento de Algorand en el mercado se entiende mejor al contrastar sus métricas operativas con los líderes históricos del sector:

\begingroup
\centering
\begin{longtable}{lp{3cm}p{3cm}p{3cm}}
\caption{Comparativa técnica de redes blockchain} \\
\toprule
\textbf{Criterio Técnico} & \textbf{Bitcoin (BTC)} & \textbf{Ethereum (ETH)} & \textbf{Algorand (ALGO)} \\
\midrule
\endhead
\textbf{Propósito Principal} & Reserva de valor y sistema de pagos simple. & Plataforma global de dApps y Smart Contracts. & Ecosistema financiero institucional y dApps de alta velocidad. \\
\midrule
\textbf{Mecanismo de Consenso} & \textit{Proof of Work} (PoW). Altamente intensivo en energía. & \textit{Proof of Stake} (PoS). Requiere un bloqueo de 32 ETH para validar. & \textit{Pure Proof of Stake} (PPoS). Requiere solo 1 ALGO. \\
\midrule
\textbf{Velocidad (TPS)} & ~7 transacciones/seg. & ~15-30 transacciones/seg. & +1.000 transacciones/seg. \\
\midrule
\textbf{Finalidad} & ~60 minutos (6 bloques). & ~12-15 minutos. & \textbf{< 5 segundos} (Sin forks). \\
\midrule
\textbf{Costes} & Altos y variables. & Muy variables (Gas Fees). & Fijos (\textbf{~0.001 ALGO}). \\
\midrule
\textbf{Creación de Activos} & No nativa. & ERC-20 (complejo). & Protocolo ASA (nativo). \\
\bottomrule
\end{longtable}
\endgroup

\section{Escenario de Uso: Mercado Institucional de Bonos de Carbono}

\textbf{Contexto del Problema:} El mercado voluntario de créditos de carbono corporativos sufre de ineficiencia administrativa, opacidad y riesgo de "doble conteo".

\textbf{Solución Implementada en Algorand:} Plataforma descentralizada de emisión, adquisición y liquidación en tiempo real.

\subsection{Flujo Operativo Detallado}
\begin{enumerate}
    \item \textbf{Emisión del Activo Nativo (ASA Protocol):} Se acuñan tokens "CO2-CREDIT" usando el estándar ASA, heredando la velocidad y seguridad nativa.
    \item \textbf{Despliegue del Smart Contract (Escrow Automático):} Contrato inteligente en PyTeal que actúa como depósito de garantía.
    \item \textbf{Proceso de \textit{Opt-in} y Compra:} La corporación compradora realiza un \textit{Opt-in} para aceptar el activo y luego procede a la compra.
    \item \textbf{Transferencias Atómicas (\textit{Atomic Transfers}):} El pago y la entrega se ejecutan simultáneamente. Sin riesgo de contraparte.
    \item \textbf{Quema y Trazabilidad Pública:} Al usar el bono, se envían los tokens a una \textit{Burn Address}. Registro inmutable y auditable por terceros.
\end{enumerate}

\end{document}
e}

\end{document}
